% BHU PYQ -- Manually transcribed & error-corrected
% Paper: GLB-301: Petrology and Economic Geology
% Exam: B.Sc. (Hons.) Semester III Examination, 2016-17

\documentclass[11pt,a4paper]{article}
\usepackage[a4paper,top=2.5cm,bottom=2.5cm,left=2.8cm,right=2.8cm,headheight=14pt]{geometry}
\usepackage{amsmath,amssymb}
\usepackage[T1]{fontenc}
\usepackage[utf8]{inputenc}
\usepackage{lmodern,microtype,fancyhdr,enumitem,hyperref,lastpage,parskip}

\hypersetup{colorlinks=true,urlcolor=black,linkcolor=black,
  pdftitle={GLB-301 Petrology and Economic Geology -- BHU B.Sc. Sem III 2016-17}}
\pagestyle{fancy}\fancyhf{}
\renewcommand{\headrulewidth}{0.4pt}\renewcommand{\footrulewidth}{0.4pt}
\fancyhead[L]{\small   \textit{Banaras Hindu University}}
\fancyhead[R]{\small   \textit{GLB-301: Petrology and Economic Geology}}
\fancyfoot[C]{\small Page  \thepage\ of \pageref{LastPage}}


\newlist{parts}{enumerate}{2}
\setlist[parts,1]{label=(\alph*),leftmargin=*,itemsep=5pt,topsep=3pt}

\newlist{romanparts}{enumerate}{2}
\setlist[romanparts,1]{label=(\roman*),leftmargin=*,itemsep=5pt,topsep=3pt}

\newcommand{\pts}[1]{\hfill   \textit{[#1]}}


\begin{document}

\begin{center}
  {\large\bfseries BANARAS HINDU UNIVERSITY}\\[2pt]
  {\normalsize B.Sc. (Hons.) --- Semester III Examination, 2016-17}\\[6pt]
  \rule{0.8\linewidth}{0.5pt}\\[6pt]
  {\large\bfseries Paper No. GLB-301: Petrology and Economic Geology}\\[4pt]
  {\normalsize Time: 3 Hours \hfill Maximum Marks: 70}\\[2pt]
  {\small Ref. No.: BS/Sem III/169}
\end{center}
\rule{\linewidth}{0.4pt}
\smallskip



\noindent   \textit{Instructions: Attempt   \textbf{five} questions in all, selecting   \textbf{two} questions from Section A and   \textbf{two} questions from Section B. Section C is compulsory. All questions carry equal marks (14 marks each).}
\bigskip

\bigskip


\noindent {\large\bfseries SECTION --- A}
\rule{\linewidth}{0.4pt}\smallskip




\noindent   \textbf{Question 1.} Describe the magmatic differentiation and assimilation processes. \pts{14}

\medskip


\noindent   \textbf{Question 2.} Write short notes on   \textbf{any four} of the following: \pts{$4  \times 3.5 = 14$}

\begin{parts}
  \item Cross-bedding
  \item Classification of Sandstone
  \item Dolerite
  \item Gneiss
  \item Metamorphic textures
\end{parts}

\medskip


\noindent   \textbf{Question 3.} What do you understand by clastic texture? Discuss in detail the observable properties based on which the texture of a clastic sedimentary rock is described. \pts{14}

\medskip


\noindent   \textbf{Question 4.} Briefly describe the types of metamorphism. Discuss the metamorphic Facies. \pts{14}

\bigskip


\noindent {\large\bfseries SECTION --- B}
\rule{\linewidth}{0.4pt}\smallskip




\noindent   \textbf{Question 5.} Write in detail the different classifications of ore deposits. \pts{14}

\medskip


\noindent   \textbf{Question 6.} What is coal? Discuss the stages of coal formation and the Indian distribution of bituminous coal. \pts{14}

\medskip


\noindent   \textbf{Question 7.} Write the chemical compositions, important physical properties, and uses of   \textbf{any four} of the following economic minerals: \pts{$4  \times 3.5 = 14$}

\begin{parts}
  \item Magnetite
  \item Talc
  \item Chromite
  \item Kyanite
  \item Graphite
  \item Galena
\end{parts}

\bigskip


\noindent {\large\bfseries SECTION --- C (Compulsory)}
\rule{\linewidth}{0.4pt}\smallskip




\noindent   \textbf{Question 8.} Comment on the following in two or three sentences each: \pts{$7  \times 2 = 14$}

\begin{romanparts}
  \item Write the four divisions of the hydrothermal process.
  \item Difference between ore and gangue with examples.
  \item Draw neat sketches of folded and faulted oil traps.
  \item Define metasomatism. Write the essential minerals of syenite rock.
  \item Difference between block lava (aa) and ropy lava (pahoehoe).
  \item Difference between oomicrite and oosparite.
  \item Intraformational conglomerate.
\end{romanparts}

\vfill\rule{\linewidth}{0.4pt}

\begin{center}\small
\href{https://bhu-exam.vercel.app/}{Click here to visit our website}\\[4pt]
\href{https://me-aryan.vercel.app/}{Click here to visit our contributor}\\[4pt]
\href{https://quashan-me.vercel.app/}{Click here to visit our contributor}
\end{center}
\end{document}
